\documentclass{article}
\usepackage[margin=0.8in]{geometry}
\usepackage{amsmath, amssymb, amsthm}
\usepackage{hyperref}
\usepackage{enumerate}

\newtheorem{proposition}{Proposition}
\newtheorem{definition}{Definition}
\newtheorem{theorem}{Theorem}
\newtheorem{remark}{Remark}

\title{Introduction to University Mathematics}
\author{}
\date{}

\begin{document}
\maketitle
\tableofcontents
\newpage

\section{Lecture 1}
\subsection{The natural numbers and induction}

\begin{definition}
  \normalfont
  A \underline{natural number} is a member of the sequence 0, 1, 2, ... formed by starting from 0 and successively adding 1. Set of all natural number is written as $\mathbb{N} = \{1,2,...\}$.
\end{definition}

\begin{remark}
  Inclusion of 0 in natural number is not universal. Sometimes it is not included.
\end{remark}

Natural number can be added and multiplied to produce natural numbers . i.e. If $m,n\in\mathbb{N}$ then $m+n,m*n\in\mathbb{N}$ ($m*n$ is often written as $mn$). 

0 and 1 are two special natural numbers because they are additive and multiplicative \underline{identities}. i.e. $\forall n\in\mathbb{N}, n + 0 = n \text{ and } n * 1 = n$. The natural numbers have ordering so it makes sense to have the notion comparison between two natural numbers. All the mathematical objects do not enjoy this property (for instance matrices).

\begin{definition}
  \normalfont
  Let $m, n \in\mathbb{N}$. We write $m\leq n$ to mean that $\exists k \in \mathbb{N}:m+k = n$.
\end{definition}

\begin{theorem}
  \normalfont
  (Principle of mathematical induction) Let $P(n)$ be a family of statements indexed by the natural numbers. Suppose:
  \begin{enumerate}[i.]
    \item $P(0)$ is true, and
    \item $\forall n\in\mathbb{N}$, if $P(n)$ and $P(n+1)$ is true
  \end{enumerate}
  then $P(n)$ is true for all $n\in\mathbb{N}$.
\end{theorem}

\begin{proposition}
  \normalfont
  $\forall n \in N$, \[\sum_{k=0}^{n}k = \frac{n(n+1)}{2}.\]
\end{proposition}

\begin{proof}
\end{proof}

\end{document}